\documentclass[11pt,a4paper]{article}
\usepackage[margin=0.85in]{geometry}
\usepackage[T1]{fontenc}
\usepackage{lmodern}
\usepackage{amsmath,amssymb,mathtools}
\usepackage{graphicx,booktabs,array,enumitem}
\usepackage[dvipsnames]{xcolor}
\usepackage{microtype}
\usepackage{caption}
\usepackage{listings}
\usepackage{hyperref}
\hypersetup{colorlinks=true,linkcolor=MidnightBlue,urlcolor=MidnightBlue,
  pdfauthor={},pdfsubject={Reinforcement learning worked solutions}}
\setlength{\parindent}{0pt}
\setlength{\parskip}{6pt}
\setlist{nosep,leftmargin=*}
\setlength{\emergencystretch}{2em}
\captionsetup{font=small,labelfont=bf}
\lstset{basicstyle=\ttfamily\small,breaklines=true,
  columns=fullflexible,keepspaces=true,showstringspaces=false,
  commentstyle=\color{gray},keywordstyle=\color{MidnightBlue},
  aboveskip=6pt,belowskip=6pt}
\newcommand{\E}{\mathbb{E}}
\newcommand{\R}{\mathbb{R}}
\newcommand{\norm}[1]{\left\lVert #1\right\rVert}
\newcommand{\sg}{\operatorname{stopgrad}}
\DeclareMathOperator*{\argmax}{arg\,max}
\DeclareMathOperator*{\argmin}{arg\,min}
\newcommand{\studytitle}[1]{%
  \hypersetup{pdftitle={#1 - Worked Solutions}}%
  {\small\color{MidnightBlue} BASICRL\quad /\quad WORKED SOLUTIONS}\par
  \vspace{5pt}{\LARGE\bfseries #1}\par\vspace{5pt}}
\newcommand{\historical}{\textbf{Experiment status.} The following numbers and
  curves are retained from earlier runs, not from a rerun of the revised
  reference code. They are examples for interpretation, not performance guarantees.\par}

\begin{document}
\studytitle{NumPy and PyTorch Foundations}

\section{Row-wise distances}
For $X,Y\in\R^{n\times d}$, each row is one vector, so the desired result is
\[
 D_{i,1}=\sqrt{\sum_{j=1}^d(X_{ij}-Y_{ij})^2},\qquad D\in\R^{n\times1}.
\]
The reduction must occur along the feature axis, while keeping a singleton
dimension for the result. Summing every entry instead produces one scalar.
\begin{lstlisting}
diff = x - y                       # (n, d)
summed = np.sum(diff**2, axis=1, keepdims=True)
distances = np.sqrt(summed)         # (n, 1)
\end{lstlisting}
For example, $X=[(0,0),(3,4)]$ and $Y=[(3,4),(3,0)]$ give
$D=[(5),(4)]$. Keeping the result two-dimensional prevents accidental
broadcasting of an $(n,)$ array against an $(n,1)$ target into an $(n,n)$ array.

\section{Assembling four squares}
Concatenate along columns to construct each row of squares, then along rows
to combine the two rows:
\begin{lstlisting}
top = torch.cat((square_1, square_2), dim=1)
bottom = torch.cat((square_3, square_4), dim=1)
combined = torch.cat((top, bottom), dim=0)
\end{lstlisting}
Each input is $(n,n)$, each intermediate is $(n,2n)$, and the result is
$(2n,2n)$. The default concatenation dimension is 0, which would stack all
four inputs vertically rather than construct the intended grid.

\section{From video frames to a filmstrip}
Stack $N$ frames with shape $(H,W,C)$ into $(N,H,W,C)$. To join frames
horizontally without interleaving their rows, first permute to $(H,N,W,C)$:
\begin{lstlisting}
frames = torch.as_tensor(np.stack(video_frames, axis=0))
n, h, w, c = frames.shape
filmstrip = frames.permute(1, 0, 2, 3).reshape(h, n*w, c)
\end{lstlisting}
The pixel mapping is
\[
 \mathrm{filmstrip}[h,nW+w,c]=\mathrm{frames}[n,h,w,c].
\]
A direct reshape of $(N,H,W,C)$ to $(H,NW,C)$ does not generally satisfy this
identity. The input list is assumed nonempty with equal frame shapes.

\newpage
\section{Linear regression and optimization}
For a batch $X\in\R^{B\times d}$ and targets $Y\in\R^{B\times m}$, a linear
layer computes $\widehat Y=XW^\top+b$ with $W\in\R^{m\times d}$. Train it by
minimizing mean squared error:
\[
 L(W,b)=\frac{1}{Bm}\sum_{i=1}^{B}\sum_{j=1}^{m}
 (\widehat Y_{ij}-Y_{ij})^2.
\]
\begin{lstlisting}
def forward(self, x):
    return self.linear(x)

loss_fn = torch.nn.MSELoss()
optimizer = torch.optim.Adam(model.parameters(), lr=0.01)
for epoch in range(num_epochs):
    optimizer.zero_grad()
    loss = loss_fn(model(x), y)
    loss.backward()
    optimizer.step()
\end{lstlisting}
Clearing gradients prevents accumulation across updates. Backpropagation
computes parameter gradients; the optimizer then changes the parameters.
Targets and predictions must have the same shape. In the scalar-output case,
both should be $(B,1)$.

\subsection{Checkpoint selection}
The exercise provides training and test data, without a separate validation
set. Therefore the reference implementation selects the checkpoint by
\emph{training} loss and reserves the test set for final evaluation. In a
larger experiment, introduce a validation set for model selection.

A saved parameter state must correspond to the loss used to select it.
The reference saves a new best state before the optimizer update and also
evaluates the final update. Calling \texttt{torch.save} immediately writes
the state to disk; holding \texttt{state\_dict()} in memory for later use
would require copying the tensors to preserve that earlier state.

\begin{lstlisting}
state = torch.load(path, map_location="cpu", weights_only=True)
out_dim, in_dim = state["linear.weight"].shape
model = LinearRegression(in_dim, out_dim)
model.load_state_dict(state)
model.eval()
with torch.no_grad():
    predictions = model(test_x)
\end{lstlisting}

\section{Checks that reveal real mistakes}
Verify distances on unequal rows, concatenate squares with distinct values,
and construct a filmstrip with multiple rows and frames so ordering errors
are visible. For regression, check that loss decreases, predictions retain
their target shape, and checkpoint loading reproduces the saved model's
predictions. A deterministic linear target is useful for a small smoke test;
it does not establish performance on arbitrary datasets.

\textbf{Solution scope.} The array identities are exact under the stated
shape assumptions. Regression convergence depends on the data, optimization
settings, and training duration; no universal convergence claim is made.
\end{document}
